\section{Exceptional points as multiple zeros of the pole condition}
\label{sec:exceptional-points}
% BEGIN CURATED SUBJECT INDEX
\index{boundary condition!incoming and outgoing}
\index{exceptional point}
\index{Jordan chain}
\index{self-orthogonality}
% END CURATED SUBJECT INDEX

\calculationroute
  {A parameter-dependent incoming coefficient and left--right resonance
  normalization.}
  {Expand the analytic pole condition at a multiple zero and translate the
  result into a Jordan chain and self-orthogonality.}
  {Distinguish an exceptional resonance from an accidental crossing or a
  discretized-continuum artifact.}

\subsection{Degeneracy, defectiveness, and Jordan chains}
\label{subsec:ep-definition}
% BEGIN CURATED SUBJECT INDEX
\index{operator!adjoint and self-adjoint}
\index{operator!closed}
\index{spectrum!point}
\index{resolvent}
\index{Feshbach projection}
\index{self-energy}
\index{exceptional point}
\index{Jordan chain}
\index{exceptional point!defectiveness}
% END CURATED SUBJECT INDEX

Let $H(\lambda)$ be an analytic family of closed operators or of finite-basis
approximations, with complex control parameter $\lambda$.  A spectral
degeneracy at $\lambda=\vsub{\lambda}{EP}$ is an exceptional point if the
algebraic multiplicity exceeds the geometric multiplicity.  For a second-order
exceptional point, there is one eigenvector $\ket{\chi_0}$ and one associated
vector $\ket{\chi_1}$ satisfying
\begin{align}
  \left(\vsub{H}{EP}-\vsub{E}{EP}\right)
  \ket{\chi_0}&=0,
  \notag\\
  \left(\vsub{H}{EP}-\vsub{E}{EP}\right)
  \ket{\chi_1}&=\ket{\chi_0},
  \label{eq:jordan-chain}
\end{align}
where $\vsub{H}{EP}=H(\vsub{\lambda}{EP})$.  The restriction of
$\vsub{H}{EP}$ to the root subspace is a $2\times2$ Jordan block.  A
diabolic degeneracy of a self-adjoint family, by contrast, has two independent
eigenvectors and no Jordan chain.

Near a generic second-order exceptional point, the characteristic equation
has the local form
\begin{equation}
  (E-\vsub{E}{EP})^2
  -c(\lambda-\vsub{\lambda}{EP})
  +\text{higher-order terms}=0,
  \label{eq:ep-characteristic}
\end{equation}
where $c\ne0$ is a complex coefficient.  The eigenvalues have Puiseux
expansions
\begin{equation}
  E_{\pm}(\lambda)
  =\vsub{E}{EP}
  \pm\alpha(\lambda-\vsub{\lambda}{EP})^{1/2}
  +\order(\lambda-\vsub{\lambda}{EP}),
  \label{eq:ep-puiseux}
\end{equation}
with $\alpha^2=c$.  Encircling the branch point once exchanges $E_+$ and
$E_-$; two encirclements return the eigenvalue to its original sheet.

The square root is generic but not universal.  A $p$th-order Jordan block gives
fractional powers $(\lambda-\vsub{\lambda}{EP})^{1/p}$, while a continuum
threshold can modify the exponent through the energy dependence of the
self-energy in Eq.~\eqref{eq:feshbach-H}.  The local singularity must be
derived from the continued resolvent rather than inferred only from a finite
matrix fit
\cite{Ashida2020NonHermitian,Garmon2021AnomalousEP}.

\subsection{Self-orthogonality and eigenvalue condition numbers}
\label{subsec:self-orthogonality}
% BEGIN CURATED SUBJECT INDEX
\index{spectrum!point}
\index{pseudospectrum}
\index{biorthogonality}
\index{complex-symmetric operator}
\index{exceptional point}
\index{self-orthogonality}
% END CURATED SUBJECT INDEX

For a simple eigenvalue, choose right and left eigenvectors
$\ket{\Psi^{\mathrm R}}$ and $\bra{\Psi^{\mathrm L}}$.  First-order
perturbation theory contains the denominator
\begin{equation}
  \braket{\Psi^{\mathrm L}|\Psi^{\mathrm R}} .
  \label{eq:left-right-overlap}
\end{equation}
At an exceptional point, the left eigenvector is orthogonal to the coalesced
right eigenvector in the biorthogonal pairing:
\begin{equation}
  \braket{\chi_0^{\mathrm L}|\chi_0^{\mathrm R}}=0 .
  \label{eq:ep-self-orthogonality}
\end{equation}
If one enforces unit overlap away from the EP, the individual vector norms
diverge as the EP is approached.  A useful sensitivity measure is
\begin{equation}
  \mathcal K_n
  =\frac{\lVert\Psi_n^{\mathrm R}\rVert
  \lVert\Psi_n^{\mathrm L}\rVert}
  {\left|\braket{\Psi_n^{\mathrm L}|\Psi_n^{\mathrm R}}\right|} .
  \label{eq:eigenvalue-condition-number}
\end{equation}
The condition number $\mathcal K_n$ diverges at self-orthogonality.  Large
$\mathcal K_n$ also means that a small operator perturbation can move the
eigenvalue substantially.  This links EP physics to pseudospectral
sensitivity.

For a complex-symmetric representation, the overlap in
Eq.~\eqref{eq:ep-self-orthogonality} becomes the $c$ product.  A vanishing
$c$ norm is therefore not evidence that the wave function itself vanishes; it
is the diagnostic that right and left eigendirections have coalesced.

\subsection{Exceptional points in an open scattering problem}
\label{subsec:ep-scattering}
% BEGIN CURATED SUBJECT INDEX
\index{topology}
\index{spectrum!point}
\index{Riemann sheet}
\index{pseudostate}
\index{CDCC!bin discretization}
\index{exceptional point}
\index{exceptional point!defectiveness}
% END CURATED SUBJECT INDEX

There are three distinct situations that are sometimes grouped together.
First, two isolated resonance poles may coalesce on a nonphysical sheet.
Second, a resonance pole may meet a threshold branch point.  Third, a pole may
meet the complex-scaled continuum in a particular $L^2$ representation.  Only
when the analytic family develops a defective root subspace is the event an
exceptional point in the operator-theoretic sense.  A visual crossing of a
finite-basis continuum ray is not sufficient.

The complex-scaling method provides a controlled setting for the third
situation.  An exposed resonance is a discrete eigenvalue of $H_\theta$.  As a
physical parameter $\lambda$ changes, its pole trajectory can reach the
rotated continuum.  In a momentum-bin discretization, one can identify a
particular continuum pseudostate that coalesces with the resonance, verify the
vanishing discriminant, and follow the resulting branch exchange.  The
continuum limit must then be taken through smeared states or a specified test
space
\cite{Morikawa:2025inx}.

Changing the scaling angle $\theta$ at fixed physical Hamiltonian is different.
The pole does not move in the exact problem; only the boundary of its $L^2$
representation moves.  Loss of exposure at a critical angle need not be a
physical EP.  It becomes an EP problem only if the chosen analytic family and
root subspace satisfy the defectiveness conditions.  This distinction avoids
assigning topological charge to a numerical continuum crossing created solely
by truncation.

\subsection{Feshbach description of pole coalescence}
\label{subsec:ep-feshbach}
% BEGIN CURATED SUBJECT INDEX
\index{spectrum!point}
\index{resolvent}
\index{resonance!residue}
\index{Feshbach projection}
\index{effective Hamiltonian}
\index{exceptional point}
% END CURATED SUBJECT INDEX

In a two-dimensional $P$ space, the effective Hamiltonian can be written
\begin{equation}
  \vsub{H}{eff}(z,\lambda)
  =H_{PP}(\lambda)+\Sigma(z,\lambda) .
  \label{eq:ep-effective-H}
\end{equation}
The pole equation is
\begin{equation}
  D(z,\lambda)
  =\det\left[z-\vsub{H}{eff}(z,\lambda)\right]=0 .
  \label{eq:ep-determinant}
\end{equation}
A second-order pole coalescence satisfies
\begin{equation}
  D(\vsub{z}{EP},\vsub{\lambda}{EP})=0,
  \qquad
  \left.\frac{\partial D}{\partial z}\right|_{\mathrm{EP}}=0 .
  \label{eq:ep-double-root}
\end{equation}
Generically, a second control condition is needed if all parameters are
restricted to be real.  Complex continuation reduces the local geometry to
the square-root form in Eq.~\eqref{eq:ep-puiseux}.  If
$\partial_z\Sigma$ diverges at a threshold, the balance in
Eq.~\eqref{eq:ep-double-root} changes and noninteger orders other than $1/2$
can appear.

This formulation makes clear that an EP is a statement about both eigenvalues
and residues.  At a double pole, the resolvent contains
\begin{equation}
  \Resolvent(z)
  =\frac{A_{-2}}{(z-\vsub{z}{EP})^2}
  +\frac{A_{-1}}{z-\vsub{z}{EP}}
  +\vsub{\Resolvent}{reg}(z) ,
  \label{eq:ep-resolvent-laurent}
\end{equation}
where $A_{-2}$ and $A_{-1}$ act on the Jordan root subspace.  A fit of pole
locations without residues cannot distinguish all possible near-degenerate
line shapes.

\subsection{Dynamics around an exceptional point}
\label{subsec:ep-dynamics}
% BEGIN CURATED SUBJECT INDEX
\index{monodromy}
\index{adiabatic approximation}
\index{exceptional point}
\index{geometric phase}
% END CURATED SUBJECT INDEX

The spectral monodromy of Eq.~\eqref{eq:ep-puiseux} is independent of how a
parameter loop is traversed.  Actual time-dependent encircling is more subtle.
The adiabatic theorem for a non-normal decaying system can fail because decay
rates select the least lossy branch and exponentially amplify small
nonadiabatic components.  A geometric phase extracted from analytic
continuation should therefore be distinguished from a state-transfer protocol
in real time.  Experiments on dissipative qubits and monitored systems probe
both aspects but require an explicit unraveling or measurement model
\cite{Foroozani:2015jfz,Naghiloo:2019dqw,
Breuer2016NonMarkovian,Gneiting:2020pew}.
